\seccionanexo{Anexo A: Cuestionarios de Diagnóstico del Proceso de Inscripción (Objetivo Específico 1)}

\providecommand{\casilla}{\framebox[0.38cm]{\rule{0pt}{0.33cm}}}

Cuestionarios de preguntas cerradas con diseño de respuesta para diagnosticar las deficiencias operativas y los requerimientos funcionales del proceso de inscripción y matriculación del BTH, conforme al objetivo específico 1 y al diagnóstico 2022-2025 del Capítulo 2. Cada bloque presenta una cuadrícula de casillas que el encuestado marca con una X. Escala Likert de 5 puntos: 1 = Totalmente en desacuerdo, 2 = En desacuerdo, 3 = Neutral, 4 = De acuerdo, 5 = Totalmente de acuerdo.

\noindent\textbf{Cuestionario D-1: Diagnóstico del proceso de inscripción --- Estudiantes postulantes (n = 60, autoadministrado)}

\begin{table}[H]
    \centering
    \caption{Encabezado e instrucciones del cuestionario D-1 (estudiantes).}
    \label{tab:d1_encabezado}
    \begin{tabularx}{\textwidth}{p{0.26\textwidth} X}
        \toprule
        \textbf{Campo} & \textbf{Contenido} \\
        \midrule
        Institución & MÓDULO TECNOLÓGICO PRODUCTIVO SAN JULIÁN BTH --- Especialidad de Sistemas Informáticos \\
        Instrumento & Cuestionario D-1: Diagnóstico del proceso de inscripción y matriculación (OE1) \\
        Muestra & 60 estudiantes postulantes de las siete especialidades (muestreo por conveniencia) \\
        Aplicación & Primer bimestre de la gestión 2026 \\
        Tiempo estimado & 6 minutos \\
        Instrucciones & Marcar con una X una sola opción por pregunta. Responder según lo vivido en los procesos 2022-2025. Respuestas anónimas (C.I. como ****). \\
        \bottomrule
    \end{tabularx}
\end{table}

\Needspace{7cm}
\noindent\textbf{Bloque A. Identificación (anonimizada):}

\begin{table}[H]
    \centering
    \caption{Diseño de respuesta del bloque A (identificación) del cuestionario D-1.}
    \label{tab:d1_identificacion}
    \renewcommand{\arraystretch}{1.60}
    \begin{tabularx}{\textwidth}{c L}
        \toprule
        \textbf{N.º} & \textbf{Pregunta y opciones de respuesta} \\
        \midrule
        A1 & Unidad educativa de procedencia: \rule[-2pt]{4.5cm}{0.4pt} \\
        \addlinespace
        A2 & Especialidad de interés: \mbox{\casilla{} Construcción Civil}\quad \mbox{\casilla{} Contaduría General}\quad \mbox{\casilla{} Electrónica}\quad \mbox{\casilla{} Gastronomía}\quad \mbox{\casilla{} Mecánica Automotriz}\quad \mbox{\casilla{} Sistemas Informáticos}\quad \mbox{\casilla{} Textiles y Confecciones} \\
        \addlinespace
        A3 & Edad: \mbox{\casilla{} 15 años o menos}\quad \mbox{\casilla{} 16 años}\quad \mbox{\casilla{} 17 años}\quad \mbox{\casilla{} 18 años o más} \\
        \addlinespace
        A4 & Acceso a internet en su teléfono celular: \mbox{\casilla{} Siempre}\quad \mbox{\casilla{} A veces}\quad \mbox{\casilla{} Nunca} \\
        \bottomrule
    \end{tabularx}
\end{table}

\Needspace{15cm}
\noindent\textbf{Bloque B. Preguntas cerradas sobre la experiencia de inscripción (2022-2025):}

\begin{table}[H]
    \centering
    \caption{Diseño de respuesta del bloque B (preguntas cerradas) del cuestionario D-1.}
    \label{tab:d1_cerradas}
    \renewcommand{\arraystretch}{1.60}
    \begin{tabularx}{\textwidth}{c L}
        \toprule
        \textbf{N.º} & \textbf{Pregunta y opciones de respuesta} \\
        \midrule
        B1 & ¿Participó en inscripciones presenciales del MTP antes de la gestión 2026?\newline
            \mbox{\casilla{} Sí}\quad \mbox{\casilla{} No} \\
        \addlinespace
        B2 & ¿Cuánto sumó su tiempo de espera en fila en la inscripción presencial?\newline
            \mbox{\casilla{} Ninguno}\quad \mbox{\casilla{} Menos de 1 día}\quad \mbox{\casilla{} 1 a 2 días}\quad \mbox{\casilla{} 3 a 5 días}\quad \mbox{\casilla{} Más de 5 días} \\
        \addlinespace
        B3 & ¿Cuánto gastó en pasajes y fotocopias durante el proceso?\newline
            \mbox{\casilla{} Menos de 20}\quad \mbox{\casilla{} 20 a 50}\quad \mbox{\casilla{} 51 a 100}\quad \mbox{\casilla{} Más de 100 Bs} \\
        \addlinespace
        B4 & ¿Olvidó su contraseña de la página de preinscripción de la gestión 2025?\newline
            \mbox{\casilla{} Sí}\quad \mbox{\casilla{} No} \\
        \addlinespace
        B5 & ¿Cómo resolvió la recuperación de su contraseña?\newline
            \mbox{\casilla{} La recuperé por internet sin ayuda}\quad \mbox{\casilla{} Acudí al Módulo y esperé varios días}\quad \mbox{\casilla{} No pude recuperarla}\quad \mbox{\casilla{} No la olvidé} \\
        \addlinespace
        B6 & ¿Qué le ocasionó principalmente esa dificultad?\newline
            \mbox{\casilla{} Pérdida de tiempo}\quad \mbox{\casilla{} Pérdida del cupo que deseaba}\quad \mbox{\casilla{} Gasto adicional}\quad \mbox{\casilla{} Ninguna} \\
        \addlinespace
        B7 & ¿Qué documentos exigía el BTH para completar la matriculación (libretas de 1ro a 3ro, cédula del tutor, carné y certificado de nacimiento)?\newline
            \mbox{\casilla{} Ninguno}\quad \mbox{\casilla{} Algunos}\quad \mbox{\casilla{} Todos} \\
        \addlinespace
        B8 & ¿En qué estado quedó su expediente en la gestión 2025?\newline
            \mbox{\casilla{} Completo}\quad \mbox{\casilla{} Incompleto}\quad \mbox{\casilla{} Extravié documentos}\quad \mbox{\casilla{} No sé} \\
        \addlinespace
        B9 & ¿Sabía cómo se asignaba su curso y paralelo (4to A o 4to B)?\newline
            \mbox{\casilla{} Sí, conocía los criterios}\quad \mbox{\casilla{} Parcialmente}\quad \mbox{\casilla{} No} \\
        \addlinespace
        B10 & ¿Presentó conflictos de horario entre el Módulo y su unidad educativa?\newline
            \mbox{\casilla{} Sí}\quad \mbox{\casilla{} No} \\
        \bottomrule
    \end{tabularx}
\end{table}

\Needspace{22cm}
\noindent\textbf{Bloque C. Ítems diagnósticos en escala Likert (1--5):}

\begin{table}[H]
    \centering
    \caption{Cuadrícula de respuesta del bloque C (escala Likert 1--5) del cuestionario D-1.}
    \label{tab:d1_likert}
    \renewcommand{\arraystretch}{1.15}
    \small
    \begin{tabularx}{\textwidth}{X c c c c c}
        \toprule
        \textbf{Ítem} & \textbf{1} & \textbf{2} & \textbf{3} & \textbf{4} & \textbf{5} \\
        \midrule
        C1. Durante el proceso de inscripción presencial tuve que hacer filas de varios días o permanecer de noche en el Módulo. & \casilla & \casilla & \casilla & \casilla & \casilla \\
        C2. La espera a la intemperie (sol, lluvia o frío) me obligó a dejar de asistir al colegio o a mis actividades de trabajo. & \casilla & \casilla & \casilla & \casilla & \casilla \\
        C3. En la gestión 2025 olvidé mi contraseña de la página de preinscripción. & \casilla & \casilla & \casilla & \casilla & \casilla \\
        C4. Para recuperar mi contraseña tuve que trasladarme presencialmente al Módulo y esperar varios días. & \casilla & \casilla & \casilla & \casilla & \casilla \\
        C5. Mientras esperaba el restablecimiento de mi contraseña, los cupos de la especialidad que deseaba se agotaron. & \casilla & \casilla & \casilla & \casilla & \casilla \\
        C6. Considero necesario recuperar la contraseña con mi cédula de identidad y fecha de nacimiento, sin asistir al Módulo y en menos de 2 minutos. & \casilla & \casilla & \casilla & \casilla & \casilla \\
        C7. No sabía con exactitud qué documentos exigía el BTH para completar mi matriculación (libretas de 1ro a 3ro, cédula del tutor, carné y certificado de nacimiento). & \casilla & \casilla & \casilla & \casilla & \casilla \\
        C8. La entrega de mis documentos se realizó sin un registro claro y ordenado que me permitiera verificar su estado. & \casilla & \casilla & \casilla & \casilla & \casilla \\
        C9. Mi expediente quedó incompleto o se extravió algún documento durante el proceso de la gestión 2025. & \casilla & \casilla & \casilla & \casilla & \casilla \\
        C10. Considero fundamental que el sistema registre y verifique todos mis documentos y me entregue un comprobante con código de verificación. & \casilla & \casilla & \casilla & \casilla & \casilla \\
        C11. Desconocía los criterios con los que se asignaría mi curso y paralelo (4to A o 4to B) según mi unidad educativa de procedencia. & \casilla & \casilla & \casilla & \casilla & \casilla \\
        C12. La asignación manual de paralelos generó cursos con cantidades desiguales de estudiantes o errores de registro. & \casilla & \casilla & \casilla & \casilla & \casilla \\
        C13. Presenté conflictos de horario entre el Módulo y mi unidad educativa de procedencia. & \casilla & \casilla & \casilla & \casilla & \casilla \\
        C14. Espero que mi asignación a curso y paralelo responda a criterios equitativos, transparentes y compatibles con mi horario escolar. & \casilla & \casilla & \casilla & \casilla & \casilla \\
        C15. Prefiero realizar mi inscripción por internet desde mi teléfono celular, incluso con conexión de baja intensidad. & \casilla & \casilla & \casilla & \casilla & \casilla \\
        C16. La espera y los gastos de traslado (pasajes y fotocopias) afectaron económicamente a mi familia durante el proceso de inscripción. & \casilla & \casilla & \casilla & \casilla & \casilla \\
        \bottomrule
    \end{tabularx}
    \notatabla{Escala Likert: 1 Totalmente en desacuerdo, 2 En desacuerdo, 3 Neutral, 4 De acuerdo, 5 Totalmente de acuerdo.}
\end{table}

\Needspace{6cm}
\noindent\textbf{Bloque D. Priorización de necesidades (opción única):}

\begin{table}[H]
    \centering
    \caption{Diseño de respuesta del bloque D (priorización) del cuestionario D-1.}
    \label{tab:d1_priorizacion}
    \renewcommand{\arraystretch}{1.60}
    \begin{tabularx}{\textwidth}{c L}
        \toprule
        \textbf{N.º} & \textbf{Pregunta y opciones de respuesta} \\
        \midrule
        D1 & ¿Cuál fue su mayor dificultad en el proceso de inscripción?\newline
            \mbox{\casilla{} Recuperar la contraseña}\quad \mbox{\casilla{} Reunir y entregar los documentos}\quad \mbox{\casilla{} Obtener la especialidad y paralelo deseados}\quad \mbox{\casilla{} Ninguna} \\
        \addlinespace
        D2 & ¿Qué función considera más importante para la página de inscripción?\newline
            \mbox{\casilla{} Recuperación autónoma de la contraseña}\quad \mbox{\casilla{} Registro y verificación de documentos}\quad \mbox{\casilla{} Asignación automática de paralelos}\quad \mbox{\casilla{} Comprobante con código de verificación} \\
        \bottomrule
    \end{tabularx}
\end{table}

\vspace{0.4cm}
\noindent\textbf{Cuestionario D-2: Diagnóstico del proceso de inscripción --- Personal administrativo y docente (n = 5, entrevista estandarizada)}

\begin{table}[H]
    \centering
    \caption{Encabezado e instrucciones del cuestionario D-2 (funcionarios).}
    \label{tab:d2_encabezado}
    \begin{tabularx}{\textwidth}{p{0.26\textwidth} X}
        \toprule
        \textbf{Campo} & \textbf{Contenido} \\
        \midrule
        Institución & MÓDULO TECNOLÓGICO PRODUCTIVO SAN JULIÁN BTH --- Especialidad de Sistemas Informáticos \\
        Instrumento & Cuestionario D-2: Diagnóstico del proceso de inscripción y matriculación (OE1) \\
        Muestra & 5 funcionarios: Director, Secretaría Académica y 3 docentes (Tutores de Curso) \\
        Aplicación & Primer bimestre de la gestión 2026 \\
        Tiempo estimado & 6 minutos \\
        Instrucciones & El aplicador marca con una X la opción elegida por el entrevistado. Responder según la experiencia de las gestiones 2022-2025. Respuestas anónimas (C.I. como ****). \\
        \bottomrule
    \end{tabularx}
\end{table}

\Needspace{6cm}
\noindent\textbf{Bloque A. Identificación (anonimizada):}

\begin{table}[H]
    \centering
    \caption{Diseño de respuesta del bloque A (identificación) del cuestionario D-2.}
    \label{tab:d2_identificacion}
    \renewcommand{\arraystretch}{1.60}
    \begin{tabularx}{\textwidth}{c L}
        \toprule
        \textbf{N.º} & \textbf{Pregunta y opciones de respuesta} \\
        \midrule
        A1 & Cargo:\newline
            \mbox{\casilla{} Director}\quad \mbox{\casilla{} Secretaría Académica}\quad \mbox{\casilla{} Docente (Tutor de Curso)} \\
        \addlinespace
        A2 & Antigüedad en el MTP:\newline
            \mbox{\casilla{} Menos de 5 años}\quad \mbox{\casilla{} 5 a 10 años}\quad \mbox{\casilla{} 11 a 15 años}\quad \mbox{\casilla{} Más de 15 años} \\
        \bottomrule
    \end{tabularx}
\end{table}

\Needspace{22cm}
\noindent\textbf{Bloque B. Preguntas cerradas sobre las deficiencias del proceso (2022-2025):}

\begin{table}[H]
    \centering
    \caption{Diseño de respuesta del bloque B (preguntas cerradas) del cuestionario D-2.}
    \label{tab:d2_cerradas}
    \renewcommand{\arraystretch}{1.60}
    \begin{tabularx}{\textwidth}{c L}
        \toprule
        \textbf{N.º} & \textbf{Pregunta y opciones de respuesta} \\
        \midrule
        B1 & ¿Cuántos restablecimientos manuales de contraseña se atendieron por gestión en ventanilla?\newline
            \mbox{\casilla{} Menos de 50}\quad \mbox{\casilla{} 50 a 100}\quad \mbox{\casilla{} 101 a 200}\quad \mbox{\casilla{} Más de 200}\quad \mbox{\casilla{} No sabe} \\
        \addlinespace
        B2 & ¿Cuántas horas semanales dedicaba su área al proceso de inscripción y matriculación?\newline
            \mbox{\casilla{} Menos de 5}\quad \mbox{\casilla{} 5 a 10}\quad \mbox{\casilla{} 11 a 20}\quad \mbox{\casilla{} Más de 20} \\
        \addlinespace
        B3 & ¿Cómo se archivaban los documentos antes de la V2.0?\newline
            \mbox{\casilla{} Carpetas físicas sin registro digital}\quad \mbox{\casilla{} Planillas impresas sin control}\quad \mbox{\casilla{} Carpeta con listado parcial}\quad \mbox{\casilla{} No sabe} \\
        \addlinespace
        B4 & ¿Qué porcentaje de expedientes quedaba incompleto o sin trazabilidad en 2025?\newline
            \mbox{\casilla{} 0 a 10}\quad \mbox{\casilla{} 11 a 30}\quad \mbox{\casilla{} 31 a 50}\quad \mbox{\casilla{} Más de 50}\quad \mbox{\casilla{} No sabe} \\
        \addlinespace
        B5 & ¿Qué errores de transcripción se observaban en la asignación manual de paralelos?\newline
            \mbox{\casilla{} Ninguno}\quad \mbox{\casilla{} Pocos}\quad \mbox{\casilla{} Frecuentes}\quad \mbox{\casilla{} Muy frecuentes} \\
        \addlinespace
        B6 & ¿Qué conflictos de horario surgían entre el MTP y las unidades educativas de origen?\newline
            \mbox{\casilla{} Ninguno}\quad \mbox{\casilla{} Pocos}\quad \mbox{\casilla{} Frecuentes}\quad \mbox{\casilla{} Muy frecuentes} \\
        \addlinespace
        B7 & ¿Qué porcentaje de postulantes terminó 2025 en una especialidad distinta a la deseada?\newline
            \mbox{\casilla{} 0 a 10}\quad \mbox{\casilla{} 11 a 30}\quad \mbox{\casilla{} 31 a 50}\quad \mbox{\casilla{} Más de 50}\quad \mbox{\casilla{} No sabe} \\
        \addlinespace
        B8 & ¿Cuál es el requisito documental imprescindible para la matriculación?\newline
            \mbox{\casilla{} Libretas de 1ro a 3ro}\quad \mbox{\casilla{} Cédula del tutor}\quad \mbox{\casilla{} Carné del estudiante}\quad \mbox{\casilla{} Certificado de nacimiento}\quad \mbox{\casilla{} Todos en conjunto} \\
        \addlinespace
        B9 & ¿Qué criterio debe ser prioritario en la asignación automática de paralelos?\newline
            \mbox{\casilla{} Misma cantidad de estudiantes por paralelo}\quad \mbox{\casilla{} Mismo colegio agrupado por paralelo}\quad \mbox{\casilla{} Compatibilidad de horarios}\quad \mbox{\casilla{} No sabe} \\
        \bottomrule
    \end{tabularx}
\end{table}

\Needspace{8cm}
\noindent\textbf{Bloque C. Grado de afectación de las brechas diagnosticadas en escala Likert (1--5):}

\begin{table}[H]
    \centering
    \caption{Cuadrícula de respuesta del bloque C (escala Likert 1--5) del cuestionario D-2.}
    \label{tab:d2_likert}
    \renewcommand{\arraystretch}{1.30}
    \begin{tabularx}{\textwidth}{X c c c c c}
        \toprule
        \textbf{Brecha} & \textbf{1} & \textbf{2} & \textbf{3} & \textbf{4} & \textbf{5} \\
        \midrule
        C1. Restablecimiento manual de contraseñas en ventanilla. & \casilla & \casilla & \casilla & \casilla & \casilla \\
        C2. Expedientes incompletos o sin registro verificable. & \casilla & \casilla & \casilla & \casilla & \casilla \\
        C3. Asignación manual de cursos y paralelos en planillas externas. & \casilla & \casilla & \casilla & \casilla & \casilla \\
        C4. Conflictos de horario entre el Módulo y los colegios. & \casilla & \casilla & \casilla & \casilla & \casilla \\
        C5. Sobrecarga de trabajo del personal administrativo y docente. & \casilla & \casilla & \casilla & \casilla & \casilla \\
        \bottomrule
    \end{tabularx}
    \notatabla{Escala Likert de afectación: 1 Nula, 2 Baja, 3 Media, 4 Alta, 5 Crítica.}
\end{table}

\notatabla{Cuestionarios de diagnóstico del objetivo específico 1, aplicados en el primer bimestre de 2026 a 60 estudiantes postulantes y 5 funcionarios del MTP San Julián (muestreo por conveniencia). Datos anonimizados (C.I. como ****). Validación por juicio del tutor y prueba piloto con 5 estudiantes (ver Anexo B). Los ítems se alinean con las tres brechas del Capítulo 2: recuperación de credenciales, matriculación documental y distribución por unidad educativa.}
